
\documentclass[11pt]{article}
\usepackage[a4paper,margin=1in]{geometry}
\usepackage{amsmath,amssymb,amsthm,mathtools}
\usepackage{enumitem}
\usepackage{hyperref}
\usepackage{booktabs}

\newtheorem{theorem}{Theorem}
\newtheorem{lemma}{Lemma}
\newtheorem{proposition}{Proposition}
\newtheorem{definition}{Definition}
\newtheorem{corollary}{Corollary}
\newtheorem{warning}{Warning}
\newtheorem{example}{Example}

\newcommand{\K}{\mathsf K}
\newcommand{\C}{\mathsf C}
\newcommand{\AL}{\mathsf{AL}}
\newcommand{\Ran}{\operatorname{Ran}}
\newcommand{\Ker}{\operatorname{Ker}}
\newcommand{\Capop}{\operatorname{Cap}}
\newcommand{\Cost}{\operatorname{Cost}}

\title{Step 48: Native Probe Adequacy and Exhaustivity for Anti-Localization Membranes}
\author{Working note}
\date{}

\begin{document}
\maketitle

\section{Purpose}

The formed-layer membrane theorem says that a formed closure has no unresolved
\emph{native} predictive recombination needles when its transported native
currency matrices satisfy
\[
        \widehat K_j \preceq \Theta .
\]
This is not automatically the same as saying that the layer has no
\emph{layer-dissolving} needles.  A declared native probe family can have blind
spots.  This note states the adequacy/exhaustivity condition needed to promote
a native membrane claim to a layer-dissolving membrane claim.

The guiding slogan is:
\[
\boxed{\text{no native needles} \Rightarrow \text{no layer-dissolving needles}
       \text{ only under an adequacy bridge}.}
\]

\section{Native and dissolving probe families}

\begin{definition}[Packaged energy and native currency]
Let \(E\) be a finite-dimensional package coordinate space, \(C_\Gamma\succeq0\)
the packaged audit form, and
\[
        L_\Gamma:E\to Y
\]
the declared native probe family.  Under null-mode legality
\(\Ker C_\Gamma \subseteq \Ker L_\Gamma\), the native currency matrix is
\[
        K_L = L_\Gamma C_\Gamma^\dagger L_\Gamma^* .
\]
\end{definition}

\begin{definition}[Layer-dissolving probe family]
A layer-dissolving probe family is a readout
\[
        D_\Gamma:E\to Z
\]
whose responses are claimed to detect all ways of dissolving or breaking the
formed layer.  The corresponding dissolving currency is
\[
        K_D = D_\Gamma C_\Gamma^\dagger D_\Gamma^* ,
\]
provided \(\Ker C_\Gamma\subseteq \Ker D_\Gamma\).  If \(D_\Gamma\) sees a
legal null direction, then its variational capacity is infinite and the
dissolving claim fails.
\end{definition}

\begin{definition}[Exact adequacy]
The native family \(L_\Gamma\) is exactly adequate for \(D_\Gamma\) when there
exists a declared readout map
\[
        A:Y\to Z
\]
such that
\[
        D_\Gamma = A L_\Gamma .
\]
Equivalently, in finite dimension,
\[
        \Ker L_\Gamma \subseteq \Ker D_\Gamma .
\]
The equivalence is algebraic: \(D_\Gamma\) descends to a linear map on the
quotient \(E/\Ker L_\Gamma\).
\end{definition}

\begin{definition}[Defective adequacy]
A defective adequacy bridge is a decomposition
\[
        D_\Gamma = A L_\Gamma + R
\]
with declared residual readout \(R:E\to Z\).  The residual currency is
\[
        K_R = R C_\Gamma^\dagger R^*
\]
under null-mode legality for \(R\).
\end{definition}

\section{Adequacy transfer}

\begin{theorem}[Exact adequacy transfer]
Assume \(D_\Gamma=A L_\Gamma\) and null-mode legality.  Then
\[
        K_D = A K_L A^* .
\]
Consequently, if the native membrane has budget
\[
        K_L\preceq \Theta_Y ,
\]
then the dissolving family has budget
\[
        K_D\preceq A\Theta_Y A^* .
\]
\end{theorem}

\begin{proof}
By direct substitution,
\[
K_D = D_\Gamma C_\Gamma^\dagger D_\Gamma^*
    = A L_\Gamma C_\Gamma^\dagger L_\Gamma^* A^*
    = A K_L A^* .
\]
The budget implication follows by congruence.
\end{proof}

\begin{theorem}[Defective adequacy transfer]
Assume
\[
        D_\Gamma = A L_\Gamma + R
\]
with residual currency \(K_R=R C_\Gamma^\dagger R^*\).  Then for every
\(t>0\),
\[
\boxed{
        K_D
        \preceq
        (1+t)A K_L A^*
        +
        (1+t^{-1})K_R .
}
\]
Consequently, if \(K_L\preceq\Theta_Y\) and \(K_R\preceq \Xi\), then
\[
\boxed{
        K_D
        \preceq
        (1+t)A\Theta_YA^*
        +
        (1+t^{-1})\Xi .
}
\]
\end{theorem}

\begin{proof}
Let
\[
        W_L = L_\Gamma C_\Gamma^{\dagger/2},
        \qquad
        W_R = R C_\Gamma^{\dagger/2}.
\]
Then
\[
        K_D=(AW_L+W_R)(AW_L+W_R)^* .
\]
For vectors \(u,v\), the elementary inequality
\[
        (u+v)(u+v)^* \preceq (1+t)uu^*+(1+t^{-1})vv^*
\]
gives the desired bound after setting \(u=AW_L\) and \(v=W_R\).
\end{proof}

\section{Blind-spot and null-mode obstruction}

\begin{theorem}[Blind-spot witness]
If exact adequacy fails, then there exists a vector \(v\in E\) such that
\[
        L_\Gamma v=0,\qquad D_\Gamma v\ne0.
\]
Thus two package states indistinguishable by the declared native family can be
distinguished by a layer-dissolving probe.
\end{theorem}

\begin{proof}
Exact adequacy is equivalent to \(\Ker L_\Gamma\subseteq\Ker D_\Gamma\).
Failure of the inclusion gives the witness.
\end{proof}

\begin{corollary}[Unbounded blind-spot risk]
If there exists a sequence \(v_n\) such that
\[
        L_\Gamma v_n=0,\qquad
        \langle v_n,C_\Gamma v_n\rangle\to0,\qquad
        \|D_\Gamma v_n\|\not\to0,
\]
then the dissolving capacity is unbounded, even though the native family does
not see the witness.  If \(v\in\Ker C_\Gamma\) and \(D_\Gamma v\ne0\), then
the dissolving capacity is infinite.
\end{corollary}

\section{Adequacy as an anti-localization gate}

\begin{definition}[Adequate membrane record]
An accepted adequate membrane record consists of:
\[
\mathsf{Adeq}=(E,C_\Gamma,L_\Gamma,D_\Gamma,A,R,\Theta_Y,\Xi,\mathsf{Status})
\]
with:
\begin{enumerate}[label=(\alph*)]
\item formed closure and exact package;
\item null-mode legality for \(L_\Gamma,D_\Gamma,R\);
\item a declared native membrane \(K_L\preceq\Theta_Y\);
\item exact adequacy \(D_\Gamma=AL_\Gamma\), or defective adequacy
      \(D_\Gamma=AL_\Gamma+R\) with \(K_R\preceq\Xi\);
\item no-smuggling: \(D_\Gamma,A,R,\Xi\) are declared before using the
      dissolving conclusion;
\item nonclaim record if only native, not dissolving, exhaustivity is accepted.
\end{enumerate}
\end{definition}

\begin{theorem}[No layer-dissolving needles from adequate native membrane]
If an adequate membrane record is accepted and
\[
        K_L\preceq \Theta_Y,\qquad K_R\preceq\Xi,
\]
then every layer-dissolving recombination \(z\in Z\) has bounded capacity:
\[
        \Capop_C(z^*D;\Ran\Gamma)
        =
        z^*K_D z
        \le
        z^*\bigl[(1+t)A\Theta_YA^*+(1+t^{-1})\Xi\bigr]z
\]
for every \(t>0\).  In the exact case \(R=0\), this reduces to
\[
        K_D\preceq A\Theta_YA^* .
\]
\end{theorem}

\section{Countermodels}

\begin{example}[Native membrane with hidden dissolving probe]
Let
\[
        E=\mathbb R^2,\qquad C_\Gamma=I,\qquad
        L_\Gamma(x_1,x_2)=x_1,\qquad
        D_\Gamma(x_1,x_2)=Mx_2 .
\]
Then \(K_L=1\), but \(K_D=M^2\).  The native membrane is harmless while the
dissolving probe can have arbitrarily large capacity.
\end{example}

\begin{example}[Fake zero budget from missing null-mode legality]
Let
\[
        C_\Gamma=\begin{pmatrix}1&0\\0&0\end{pmatrix},\qquad
        L_\Gamma(x_1,x_2)=x_1,\qquad
        D_\Gamma(x_1,x_2)=x_2.
\]
The pseudoinverse expression gives \(D_\Gamma C_\Gamma^\dagger D_\Gamma^*=0\),
but the variational dissolving capacity is infinite, because \(D_\Gamma\) sees
a zero-energy direction.  Null-mode legality is therefore a required gate.
\end{example}

\begin{warning}[Scope]
This theorem does not say that every nonnative or undeclared probe is safe.  It
says that a formed closure with an adequate native family has no
layer-dissolving needles \emph{relative to the declared dissolving family and
adequacy bridge}.  If the adequacy bridge is absent, the honest status is
\(\texttt{native\_only}\), \(\texttt{support\_only}\), or
\(\texttt{failed\_adequacy}\), not accepted layer-dissolving anti-localization.
\end{warning}

\end{document}
